% GENERATED FILE. Edit canonical lessons, then run npm run generate:books.
% canonical-source-hash: fnv1a64:1fcb27aabdb2782e
% canonical-lessons: FR-C415-station, FR-C415-rame, FR-C415-carnet, FR-C415-tramway, FR-C415-trottinette, FR-R415-revision-408-415

\chapter{En ville sans voiture}
\label{ch:fr-en-ville-sans-voiture}
\hlchaptermodality{415}

\begin{chapteropening}
\textbf{By the end of this chapter:} \emph{I can get around town without a car: the stop, the train set on a line, the book of tickets, the tram and the scooter.}

Complete the last lesson of chapter 415: I can get around town without a car: the stop, the train set on a line, the book of tickets, the tram and the scooter.
\end{chapteropening}

\section[la station]{la station — the metro station, the stop; the taxi rank; the resort}

\label{lesson:FR-C415-station}



\begin{quote}
Before the new one: say the French for the trolley, then the French for the shuttle.
\end{quote}

\subsection*{You'll want to know: la station}
\textbf{la station} — \textquotedblleft{}the metro station, the stop; the taxi rank; the resort\textquotedblright{}.

\begin{quote}
Je descends à la prochaine station de métro.
\end{quote}

\begin{quote}
\emph{I get off at the next metro stop.}
\end{quote}

For trains, \emph{la gare}; for the metro, \emph{la station}. \emph{Une station de ski} — a ski resort.

\begin{cousinweb}
\emph{station} comes from Old French \emph{estation}, a learned borrowing from Latin \emph{statiō} 'a standing, a post', from \emph{stāre} 'to stand' (EtymDB).
\end{cousinweb}

\begin{grammarlens}[title={What you've built}]
One more word for telling, explaining and describing at B1.
\end{grammarlens}

\subsection*{Your turn}
\begin{itemize}
\raggedright
  \item \emph{Say it:} \emph{la station}
  \item \emph{Say it:} \emph{la station}, once more
  \item \emph{Say it:} a sentence of your own with \emph{la station}
  \item \emph{Recall:} say the French for the trolley, then the French for the shuttle, then say \emph{la station} again
\end{itemize}

\begin{tcolorbox}[breakable,colback=teal!4,colframe=teal!35!black,title={Before you move on}]
What does la station mean? (\textquotedblleft{}The metro station, the stop; the taxi rank; the resort\textquotedblright{}.) Say it once more.
\end{tcolorbox}
\section[la rame]{la rame — the metro train, the set of carriages}

\label{lesson:FR-C415-rame}



\begin{quote}
Before the new one: say the French for the shuttle, then the French for the metro station.
\end{quote}

\subsection*{You'll want to know: la rame}
\textbf{la rame} — \textquotedblleft{}the metro train, the set of carriages\textquotedblright{}.

\begin{quote}
La rame de métro est pleine à huit heures.
\end{quote}

\begin{quote}
\emph{The metro train is full at eight o'clock.}
\end{quote}

\emph{La rame} is also an oar; that is a different word.

\begin{cousinweb}
\emph{rame} in this sense comes from \emph{rame} 'ream, stack', borrowed through Catalan \emph{raima} (EtymDB) from Arabic \emph{rizma} 'bundle': a row of carriages, like a stack of sheets.
\end{cousinweb}

\begin{grammarlens}[title={What you've built}]
One more word for telling, explaining and describing at B1.
\end{grammarlens}

\subsection*{Your turn}
\begin{itemize}
\raggedright
  \item \emph{Say it:} \emph{la rame}
  \item \emph{Say it:} \emph{la rame}, once more
  \item \emph{Say it:} a sentence of your own with \emph{la rame}
  \item \emph{Recall:} say the French for the shuttle, then the French for the metro station, then say \emph{la rame} again
\end{itemize}

\begin{tcolorbox}[breakable,colback=teal!4,colframe=teal!35!black,title={Before you move on}]
What does la rame mean? (\textquotedblleft{}The metro train, the set of carriages\textquotedblright{}.) Say it once more.
\end{tcolorbox}
\section[le carnet]{le carnet — the book of tickets; the notebook}

\label{lesson:FR-C415-carnet}



\begin{quote}
Before the new one: say the French for the metro station, then the French for the metro train.
\end{quote}

\subsection*{You'll want to know: le carnet}
\textbf{le carnet} — \textquotedblleft{}the book of tickets; the notebook\textquotedblright{}.

\begin{quote}
Elle achète un carnet de dix tickets.
\end{quote}

\begin{quote}
\emph{She buys a book of ten tickets.}
\end{quote}

\emph{Un carnet} is also a small notebook.

\begin{cousinweb}
\emph{carnet} comes from Middle French \emph{quernet} 'quire of four sheets' (EtymDB), from Latin \emph{quaternī} 'four each'.
\end{cousinweb}

\begin{grammarlens}[title={What you've built}]
One more word for telling, explaining and describing at B1.
\end{grammarlens}

\subsection*{Your turn}
\begin{itemize}
\raggedright
  \item \emph{Say it:} \emph{le carnet}
  \item \emph{Say it:} \emph{le carnet}, once more
  \item \emph{Say it:} a sentence of your own with \emph{le carnet}
  \item \emph{Recall:} say the French for the metro station, then the French for the metro train, then say \emph{le carnet} again
\end{itemize}

\begin{tcolorbox}[breakable,colback=teal!4,colframe=teal!35!black,title={Before you move on}]
What does le carnet mean? (\textquotedblleft{}The book of tickets; the notebook\textquotedblright{}.) Say it once more.
\end{tcolorbox}
\section[le tramway]{le tramway — the tram}

\label{lesson:FR-C415-tramway}



\begin{quote}
Before the new one: say the French for the metro train, then the French for the book of tickets.
\end{quote}

\subsection*{You'll want to know: le tramway}
\textbf{le tramway} — \textquotedblleft{}the tram\textquotedblright{}.

\begin{quote}
Le tramway passe toutes les cinq minutes.
\end{quote}

\begin{quote}
\emph{The tram comes every five minutes.}
\end{quote}

In speech, \emph{le tram}.

\begin{cousinweb}
\emph{tramway} was borrowed from English \emph{tramway} (EtymDB). English \emph{tram} 'rail' comes from Middle Dutch \emph{trame} 'beam'.
\end{cousinweb}

\begin{grammarlens}[title={What you've built}]
One more word for telling, explaining and describing at B1.
\end{grammarlens}

\subsection*{Your turn}
\begin{itemize}
\raggedright
  \item \emph{Say it:} \emph{le tramway}
  \item \emph{Say it:} \emph{le tramway}, once more
  \item \emph{Say it:} a sentence of your own with \emph{le tramway}
  \item \emph{Recall:} say the French for the metro train, then the French for the book of tickets, then say \emph{le tramway} again
\end{itemize}

\begin{tcolorbox}[breakable,colback=teal!4,colframe=teal!35!black,title={Before you move on}]
What does le tramway mean? (\textquotedblleft{}The tram\textquotedblright{}.) Say it once more.
\end{tcolorbox}
\section[la trottinette]{la trottinette — the scooter (kick or electric)}

\label{lesson:FR-C415-trottinette}



\begin{quote}
Before the new one: say the French for the book of tickets, then the French for the tram.
\end{quote}

\subsection*{You'll want to know: la trottinette}
\textbf{la trottinette} — \textquotedblleft{}the scooter (kick or electric)\textquotedblright{}.

\begin{quote}
Mon voisin va au bureau en trottinette électrique.
\end{quote}

\begin{quote}
\emph{My neighbour goes to the office on an electric scooter.}
\end{quote}

\emph{Une trottinette électrique} — an e-scooter.

\begin{cousinweb}
\emph{trottinette} is built on \emph{trottiner} 'to trot along' with \emph{-ette} (EtymDB). \emph{Trotter} 'to trot' is usually traced to a Germanic word.
\end{cousinweb}

\begin{grammarlens}[title={What you've built}]
One more word for telling, explaining and describing at B1.
\end{grammarlens}

\subsection*{Your turn}
\begin{itemize}
\raggedright
  \item \emph{Say it:} \emph{la trottinette}
  \item \emph{Say it:} \emph{la trottinette}, once more
  \item \emph{Say it:} a sentence of your own with \emph{la trottinette}
  \item \emph{Recall:} say the French for the book of tickets, then the French for the tram, then say \emph{la trottinette} again
\end{itemize}

\begin{tcolorbox}[breakable,colback=teal!4,colframe=teal!35!black,title={Before you move on}]
What does la trottinette mean? (\textquotedblleft{}The scooter (kick or electric)\textquotedblright{}.) Say it once more.
\end{tcolorbox}
\section[(dialogue)]{Retour sur les chapitres 408–415}

\label{lesson:FR-R415-revision-408-415}



\begin{quote}
Say the French for the fine and the speed camera.
\end{quote}

\subsection*{Your turn}
Say these aloud:

\begin{itemize}
\raggedright
  \item the fine, the speed camera, to flash, to fine, to inspect
  \item the indicator, the headlight, the bend, the road surface, to turn aside
  \item the traffic jam, the gridlock, the diversion, the rest area, to overtake
  \item to crash into, to bump into, the pedestrian, the cyclist, to note
  \item the ticket office, the fare, the ride, the flat-rate pass, to validate
  \item the track, to stamp, the ticket inspector, the carriage, to stop at
  \item the berth, the end of the line, the left-luggage office, the trolley, the shuttle
  \item the metro station, the metro train, the book of tickets, the tram, the scooter
\end{itemize}

\begin{tcolorbox}[breakable,colback=teal!4,colframe=teal!35!black,title={Before you move on}]
The fine? (\textbf{la contravention}.) The speed camera? (\textbf{le radar}.) To flash? (\textbf{flasher}.) To fine? (\textbf{verbaliser}.) To inspect? (\textbf{contrôler}.) The indicator? (\textbf{le clignotant}.) The headlight? (\textbf{le phare}.) The bend? (\textbf{le virage}.) The road surface? (\textbf{la chaussée}.) To turn aside? (\textbf{dévier}.) The traffic jam? (\textbf{le bouchon}.) The gridlock? (\textbf{l'embouteillage}.) The diversion? (\textbf{la déviation}.) The rest area? (\textbf{l'aire}.) To overtake? (\textbf{doubler}.) To crash into? (\textbf{percuter}.) To bump into? (\textbf{heurter}.) The pedestrian? (\textbf{le piéton}.) The cyclist? (\textbf{le cycliste}.) To note? (\textbf{constater}.) The ticket office? (\textbf{le guichet}.) The fare? (\textbf{le tarif}.) The ride? (\textbf{le trajet}.) The flat-rate pass? (\textbf{le forfait}.) To validate? (\textbf{valider}.) The track? (\textbf{la voie}.) To stamp? (\textbf{composter}.) The ticket inspector? (\textbf{le contrôleur}.) The carriage? (\textbf{le wagon}.) To stop at? (\textbf{desservir}.) The berth? (\textbf{la couchette}.) The end of the line? (\textbf{le terminus}.) The left-luggage office? (\textbf{la consigne}.) The trolley? (\textbf{le chariot}.) The shuttle? (\textbf{la navette}.) The metro station? (\textbf{la station}.) The metro train? (\textbf{la rame}.) The book of tickets? (\textbf{le carnet}.) The tram? (\textbf{le tramway}.) The scooter? (\textbf{la trottinette}.)
\end{tcolorbox}
